\documentclass[a4paper,12pt]{article}
% add all packages
\usepackage[a4paper,bottom=35mm]{geometry}
\usepackage{amsmath}
\usepackage{amsthm}
\usepackage{mdframed}
\usepackage{textcomp} 
\usepackage{siunitx}          		
\usepackage[parfill]{parskip}    		
\usepackage{graphicx}
\usepackage{xcolor}
\usepackage{caption}
\usepackage{bbm}	
\usepackage{amssymb}
\usepackage{paralist}
\usepackage{color}
\usepackage[utf8]{inputenc}
\usepackage[english,german]{babel}
\usepackage{mathtools}
\usepackage{nicefrac}
\usepackage{wrapfig}
\usepackage[mathscr]{euscript}
\usepackage{ytableau}
\usepackage{mathtools}
\usepackage{fancyhdr}
\usepackage{url}
\usepackage{hyperref}

% definiert die Sprache
\selectlanguage{German}
\setlocalecaption{german}{figure}{Figur}
\setlocalecaption{german}{table}{Tabelle}

% add command for different dates in Swiss fashion
\newcommand{\monthwordlong}[1]{\ifcase#1 \or January\or February\or March\or April\or May\or June\or July\or August\or September\or October\or November\or December\fi}
\newcommand{\monthwordshort}[1]{\ifcase#1\or Jan\or Feb\or Mar\or Apr\or May\or Jun\or Jul\or Aug\or Sept\or Okt\or Nov\or Dez\fi} 
\newcommand{\datelong}{\number\day. \monthwordlong{\the\month} \number\year}
\newcommand{\dateshort}{\number\day. \monthwordshort{\the\month} \number\year}

% capital roman numbers
\newcommand{\uproman}[1]{\uppercase\expandafter{\romannumeral#1}}

% define color of links inside a text
\definecolor{myblue}{RGB}{0, 0, 80}
\definecolor{plotblue}{RGB}{100, 100, 255}
\hypersetup{
	colorlinks=true, 
	linkcolor=myblue,
	filecolor=magenta,      
	urlcolor=plotblue,
	citecolor=myblue,
	pdfpagemode=FullScreen,
}

% placing a figure at the top of clear page
\makeatletter
\setlength{\@fptop}{0pt}
\makeatother

% add single blankpage
\def\blankpage{\clearpage\thispagestyle{empty}\addtocounter{page}{-1}\null\clearpage}

% Lückentext erstellen
\newcommand{\lineortext}[1]{\hspace{0,1cm}\underline{\hspace{0,25cm}\hphantom{#1}\hspace{0,25cm}}\hspace{0,1cm}}
% \renewcommand{\lineortext}[1]{\hspace{0,1cm}\underline{\hspace{0,25cm}#1\hspace{0,25cm}}\hspace{0,1cm}}

% math commands
\theoremstyle{definition} % don't write everything in italic
\newtheorem{myfact}{Fact}
\newtheorem{mydef}{Definition}[section]
\newenvironment{fdef}{\begin{mdframed}\begin{mydef}}{\end{mydef}\end{mdframed}}
\newtheorem{mytheo}{Theorem}
\newenvironment{ftheo}{\begin{mdframed}\begin{mytheo}}{\end{mytheo}\end{mdframed}}
\newtheorem{myprop}{Proposition}
\newtheorem{mylemma}{Lemma}[section]
\newtheorem{myaux}{Auxiliary Lemma}
\newtheorem{mycor}{Corollary}
\newtheorem{myre}{Remark}[section]
\newtheorem{myex}{Example}
\newtheorem{mynot}{Notation}

\newtheorem{satz}{Satz}[section]
\newenvironment{fsatz}{\begin{mdframed}\begin{satz}}{\end{satz}\end{mdframed}}
\newtheorem{aufg}{Aufgabe}
\newtheorem{bsp}{Beispiel}

\newcommand{\R}{\mathbb{R}}   
\newcommand{\N}{\mathbb{N}}
\newcommand{\Z}{\mathbb{Z}}
\newcommand{\Q}{\mathbb{Q}}
\newcommand{\C}{\mathbb{C}}
\newcommand{\Lim}[1]{\raisebox{0.5ex}{\scalebox{0.8}{$\displaystyle \lim_{#1}\;$}}}
\newcommand{\vvec}[2]{\left(\begin{array}{c} #1 \\ #2 \end{array}\right)}               % two dimensional vector
\newcommand{\vvector}[3]{\left(\begin{array}{c} #1 \\ #2 \\ #3 \end{array}\right)}      % three dimensional vector
\newcommand{\irow}[1]{\begin{smallmatrix}(#1)\end{smallmatrix}}                         % inline row vector
\newcommand{\icol}[1]{\left(\begin{smallmatrix}#1\end{smallmatrix}\right)}              % inline column vector
\DeclarePairedDelimiter\ceil{\lceil}{\rceil}                                            % ceil- function
\DeclarePairedDelimiter\floor{\lfloor}{\rfloor}                                         % floor - function


% header and footers
\pagestyle{fancy}
\fancyhf{}
\fancyhead[L]{}
% \fancyhead[C]{header}
\fancyhead[R]{}
\fancyfoot[R]{\thepage}
% check for more information: https://de.overleaf.com/learn/latex/Headers_and_footers


% AdminCommands
\newcommand{\name}{Jonas Guler}
\newcommand{\examdate}{26. Okt. 2022}
\newcommand{\topic}{Vektorgeometrie Repetitionsprüfung \uproman{1}}
\newcommand{\subject}{Mathematik}
\newcommand{\class}{KSH - 3GaLaWa}
\newcommand{\school}{KSH}

\begin{document}
\thispagestyle{fancy}
\lhead{\sf \large \topic \\ \small \class, \name, \examdate}
\rhead{\sf  \\  Name: \hspace*{4cm}}
\begin{small}
	\begin{center}
		\boxed{\textbf{Wichtige Informationen:}}  \\
		\vspace{2pt}
		Der Rechenweg muss vollständig und nachvollziehbar sein.\\
		Die Schlussresultate sind auf zwei Nachkommastellen zu runden oder vollständig gekürzt anzugeben.\\
		Lösen Sie die Prüfung mit einen schwarzen oder blauen Kugelschreiber. Resultate in Bleistift, roter oder grüner Farbe werden nicht berücksichtigt.\\
		Schreiben Sie auf \textbf{jedes} Blatt (inklusive Aufgabenblatt) Ihren Namen.\\
		Es sind keine Hilfsmittel erlaubt.\\
		
		! VIEL ERFOLG !
	\end{center}
\end{small}

\Aufgabe{1 (2+4)}
\begin{itemize}
    \item[a)] Vereinfache die folgende Vektorgleichung so weit wie möglich.
	\[ \frac{1}{3}\vec{a}-\left(\frac{2}{4}\vec{b}-\vec{c}-\left(\frac{1}{4}\vec{b}-\vec{c}\right)+\frac{1}{3}\vec{a}\right) + \frac{1}{4}\vec{b} \]
	
	\item[b)] Gegeben ist ein ebenes Viereck $ABCD$. Zeige vektoriell, dass die Seitenmittelpunkte des Vierecks Ecken eines Parallelogramms sind.
	%zz: sind PQRS die Seitenmittepunkte: $\vec{PQ}=\vec{SR}$
	
\end{itemize}

\newpage
\Aufgabe{2 (6+4)}
\begin{itemize}
	\item[a)] Zerlege den Vektor $\vec{d}$ nach den Vektoren $\vec{a}$, $\vec{b}$ und $\vec{c}$.
	\[ \vec{a} = \vvector{-2}{3}{7}, \vec{b} = \vvector{1}{1}{-4}, \vec{c} = \vvector{4}{-2}{0}, \vec{d} = \vvector{2}{-12}{-20} \] % x=-4,y=-2,z=-1
	
	\item[b)] Welche Punkte auf der $x$-Achse sind von Punkt $A=(-3,12,-6)$ doppelt so weit entfernt wie vom Punkt $B=(0,6,3)$?
	% (3,0,0) und (-1,0,0)
\end{itemize}

\newpage
\Aufgabe{3 (4+4)}
\begin{itemize}
	\item[a)] Bestimme für das Dreieck $ABC$ mit Schwerpunkt $S$ die dritte Ecke $A$. \[B=(1,1,4),\:C=(6,3,0) \text{ und } S=(4,0,3)\]
	% A=(5,-4,5)
	
	\item[b)] Bilden die Punkte $A=(3,-6,5),\:B=(6,5,-4),\:C=(8,8,-2)$ und $D=(5,-3,7)$ ein Parallelogramm?
	% Ja
\end{itemize}

\end{document}